\section{Introduction}\label{sec:introduction}

The Stern--Brocot and Farey constructions organize rational numbers by
continued-fraction complexity.  A natural statistical question is what a
typical rational looks like when a fixed complexity layer is biased toward
small denominators.  In the model considered here, the complexity of the
canonical regular continued fraction
\[
 x=[0;a_1,\ldots,a_r],\qquad a_i\ge1,\quad a_r\ge2,
\]
is the digit sum \(a_1+\cdots+a_r=n\), and the weight is
\(\operatorname{den}(x)^{-s}\) for a fixed \(s>2\).  The continuant in the
denominator is nonmultiplicative across the digits.  Nevertheless, a single
partial quotient absorbs all but a tight amount of the digit sum, while the
two finite contexts on its sides converge to independent arithmetic laws.

This condensation statement is stronger than the scalar asymptotic for the
normalizing sum.  It identifies the limiting context words in total
variation, the defect of the largest partial quotient, its position in the
continued fraction, the numbers of digits on either side, and the leading
factorization of the sampled denominator.  The proof also gives the scalar
asymptotic with its exact constant.  That scalar consequence becomes the
arithmetic input for a different but closely linked question: the critical
renewal structure of the classical Fibonacci partition function.

The link is supplied by Weinstein's free generators for Fibonacci
partitions \cite{Weinstein2016}.  Their letters are reduced fractions in
\((0,1)\); a letter \(p/q\) carries multiplicative weight \(q\) and an odd
cost read from its negative continued fraction.  The cost-\(2d+1\) letter
sum is termwise identical to the regular-continuant sum at digit sum
\(d+1\).  Thus the same scalar estimate that normalizes the condensation
law determines the critical generator-cost tail.  At the unique
\(\sigma_0>2\) satisfying
\[
 \frac{\zeta(\sigma_0-1)}{\zeta(\sigma_0)}=2,
\]
the generator letters form a probability distribution with a regularly
varying tail of index \(1-\sigma_0\).  Its mean is finite and its variance is
infinite.  A second-order lattice renewal theorem then produces a fractional
power correction to the critical Fibonacci partition sum, and the same tail
places the cost in the domain of attraction of a spectrally positive stable
law.

\subsection{Main results}

For a finite word \(\mathbf w\), let \(\mathcal K(\mathbf w)\) be its
regular continuant and let \(|\mathbf w|_1\) be its digit sum.  Write
\(\mathcal W_{\rm L}\) for all finite positive words, including the empty
word, and \(\mathcal W_{\rm R}\) for the canonical regular words, again
including the empty word.  Put
\[
 \rho_s=\frac{\zeta(s-1)}{\zeta(s)}.
\]
The precise definitions occur in \Cref{sec:brocot-condensation}.  The first
result is summarized as follows.

\begin{theoremA}[Total-variation condensation]
For fixed \(s>2\), sample a canonical continued fraction of digit sum \(n\)
with probability proportional to the inverse \(s\)-th power of its
denominator.  With probability tending to one it has a unique digit larger
than \(n/2\).  If \(U_n,V_n\) are the words to the left and right of that
digit, then
\[
 (U_n,V_n)\xrightarrow{\rm TV}(U,V),
\]
where \(U,V\) are independent and
\[
 \Pr\{U=\mathbf u\}=\frac{\mathcal K(\mathbf u)^{-s}}{2\rho_s},
 \qquad
 \Pr\{V=\mathbf v\}=\frac{\mathcal K(\mathbf v)^{-s}}{\rho_s}.
\]
The defect, the location data, and the normalized denominator converge as
stated in \Cref{thm:brocot-condensation}.
\end{theoremA}

The asymmetry between the two context laws is caused only by the canonical
terminal-digit convention.  A regular rational has two positive
continued-fraction expansions, whereas the canonical form retains the one
whose final digit is at least two.  The empty word and the one-letter word
\((1)\) are endpoint exceptions.  Accounting for those exceptions is also
what resolves the constant in the earlier scalar calculation.

Let
\[
 Z_n(s)=\sum_{\substack{a_1+\cdots+a_r=n\\a_i\ge1,\ a_r\ge2}}
       \mathcal K(a_1,\ldots,a_r)^{-s}.
\]
Dushistova's notation is
\(\sigma_\beta(n)=Z_n(2\beta)\).  The earlier paper obtains a fuller
expansion whose displayed leading coefficient is
\(R_s+2R_s^2\), where \(R_s=\rho_s\)
\cite{Dushistova2007}.  The exact endpoint decomposition gives the following
result.

\begin{theoremB}[Context constant and denominator layer]
For every fixed \(s>2\),
\[
 Z_n(s)\sim2\rho_s^2n^{-s}.
\]
Under the termwise change from regular to negative continued fractions, the
cost-\(2d+1\) denominator layer satisfies
\[
 b_{2d+1}(s)=Z_{d+1}(s)\sim2\rho_s^2d^{-s}.
\]
The constant \(2R_s^2\) corrects the leading coefficient in Dushistova's
expansion.  No conclusion is drawn here about its other coefficients.
\end{theoremB}

The proof is deliberately local.  The left context of digit sum zero
contributes \(R_s\), the context of digit sum one contributes another
\(R_s\), and the contexts of larger digit sum contribute
\(2(R_s-1)R_s\).  Their sum is \(2R_s^2\).  In the printed convolution the
condition excluding the empty left context is lost after the \(u=1\) term is
separated; because the remaining convolution is doubled, that empty context
is counted twice.  Restoring one subtraction removes precisely the excess
\(R_s\).  This is not a change of continued-fraction convention or exponent.

For the probabilistic theorem, the central estimate is a summable
two-context domination.  Once a digit \(a>n/2\) is singled out,
multilinearity of the continuant gives
\[
 \mathcal K(\mathbf u,a,\mathbf v)
 =a\mathcal K(\mathbf u)\mathcal K(\mathbf v)+O_{\mathbf u,\mathbf v}(1).
\]
The inverse-continuant mass of the two contexts is summable for \(s>2\).
Words with no digit above \(n/2\) are negligible: a greedy block
decomposition handles words made only of moderate digits, and a direct cut
handles the cases with one or at least two larger digits.  Dominated
convergence gives the product context law, and the discrete Scheff\'e
argument upgrades pointwise convergence to total variation.

The second part of the paper transfers only the scalar layer estimate, not
the total-variation theorem.  Let \(F_0=0,F_1=1\), and let \(R(N)\) count
representations of \(N\) as a sum of distinct members of
\(F_2,F_3,\ldots\).  On the standard Fibonacci layer
\[
 \mathcal I_m=[F_{m+1}-1,F_{m+2}-1),
\]
put
\[
 Z_m^R(t)=\sum_{N\in\mathcal I_m}R(N)^t.
\]
The free-generator parameterization yields a renewal sequence \(u_j(s)\)
whose exact layer identity is
\[
 Z_m^R(-s)=2\sum_{j=0}^{m-1}u_j(s)+u_m(s)-1.
\]
This identity and its proof are given in
\Cref{prop:weighted-generator-renewal}.

At \(s=\sigma_0\), let \(C\) be the generator-letter cost.  Set
\(\alpha=\sigma_0-1\in(1,2)\), and let \(\mu_C=\EE C\).  The arithmetic
constant from the two context masses is \(b_C=8\), while
\[
 K_C=\frac{2^{\sigma_0-1}}{\sigma_0-1}b_C.
\]

\begin{theoremC}[Critical renewal and stable law]
The critical letter law satisfies
\[
 \Pr\{C>x\}\sim K_Cx^{1-\sigma_0}.
\]
The associated renewal sequence has the second-order expansion
\[
 u_j(\sigma_0)=\frac1{\mu_C}
 +\frac{K_C}{\mu_C^2(\sigma_0-2)}j^{2-\sigma_0}
 +o(j^{2-\sigma_0}),
\]
and hence
\[
 Z_m^R(-\sigma_0)=\frac{2m}{\mu_C}
 +\frac{2K_C}{\mu_C^2(\sigma_0-2)(3-\sigma_0)}
  m^{3-\sigma_0}+o(m^{3-\sigma_0}).
\]
Moreover, centered sums of independent copies of \(C\), normalized by a
sequence \(a_n\) with \(n\Pr\{C>a_n\}\to1\), converge to the centered
spectrally positive \(\alpha\)-stable law specified in
\Cref{thm:arithmetic-critical-stable-renewal}.  One may take explicitly
\[
 a_n=\left(\frac{2^{\sigma_0+2}}{\sigma_0-1}n\right)^{1/(\sigma_0-1)}.
\]
\end{theoremC}

The noninteger power is a finite-size signature of the arithmetic critical
tail.  Regular variation is not assumed in the renewal step: it is proved
first from the denominator-layer asymptotic.  The second-order renewal
theorem of Omey and Van Gulck \cite{OmeyVanGulck2015} then applies because
the law is aperiodic, supported on the positive integers, and has finite
mean.  The stable conclusion follows from the standard one-sided
domain-of-attraction criterion \cite{Feller1971}.

\subsection{History and contribution boundary}

Moshchevitin and Zhigljavsky studied entropy sums associated with partitions
of the unit interval generated by the Farey tree
\cite{MoshchevitinZhigljavsky2004}.  Dushistova developed precise
fixed-digit-sum denominator estimates for Brocot sequences
\cite{Dushistova2007}.  The present scalar sum is exactly her local sum after
the exponent change \(s=2\beta\), and the relation is termwise rather than
asymptotic.  The contribution here is the joint total-variation structure,
the endpoint calculation of the leading constant, and the shorter scalar
proof that does not assume an auxiliary tail estimate.  The correction is
confined to the leading coefficient: the arguments here neither verify nor
reject the later coefficients in the published expansion.

Condensation under a large-sum conditioning is familiar for independent
heavy-tailed variables and for probabilistic combinatorial structures.  In
particular, total-variation one-big-jump results and Gibbs-partition
condensation provide useful comparisons
\cite{ArmendarizLoulakis2011,Stufler2020}.  They do not imply the theorem
here.  A continuant weight does not factor into weights of its digits, so the
model is not an independent allocation or a multiplicative Gibbs partition.
The proof instead exploits the exact two-sided continuant inequality and the
summability of all finite contexts.  Accordingly, the result is presented as
a structural theorem for this nonmultiplicative arithmetic family, not as a
general condensation principle.

Fibonacci representations have a long history, beginning with the canonical
representation theorem and continuing through exact formulas and extremal
questions \cite{Zeckendorf1972,Carlitz1968,Berstel2001}.  Weinstein's orbit
and generator theory gives the coordinate system used here
\cite{Weinstein2016}.  We do not reproduce that classification.  We state
the precise pieces of it needed for the weighted layer count, check the
Fibonacci indexing and endpoint convention, and derive the renewal identity.
The critical tail constant, the second-order renewal term, and the stable
limit are then consequences of the continuant calculation in this paper.

A companion manuscript by the same authors, \emph{Finite-window Zeckendorf
thermodynamics}, studies an exact two-layer finite-window correspondence,
critical coexistence, and large deviations across the pressure corner.  No
result from that manuscript is used here.  Conversely, its leading critical
renewal law requires only the finite-mean renewal theorem; the exact tail and
fractional correction proved here are not prerequisites for its principal
finite-window results.

\subsection{Organization}

\Cref{sec:brocot-condensation} defines the denominator-weighted Brocot
model, proves the context estimates and the total-variation theorem, and
derives the corrected denominator-layer asymptotic.
\Cref{sec:fibonacci-renewal} introduces the classical Fibonacci partition
function, the minimal Weinstein generator coordinates, and the exact
one-layer weighted renewal identity.  The arithmetic critical tail,
second-order renewal expansion, stable law, and Fibonacci finite-size
consequence are proved in \Cref{sec:arithmetic-critical}.  The final section
records the scope of the two mechanisms and the open directions left by the
argument.
